\section{Proponowana metoda wstępnego przetwarzania danych}

Pierwszym etapem przetwarzania będzie rozdzielenie zbioru danych na macierz cech wejściowych $X$ oraz wektor etykiet klas $y$. Do macierzy $X$ trafi 13 cech opisujących próbkę wina, natomiast do wektora $y$ przypisana zostanie informacja o klasie.

Ponieważ wszystkie cechy wejściowe są liczbowe, nie ma potrzeby kodowania danych nienumerycznych.

Najważniejszym etapem wstępnego przetwarzania będzie skalowanie danych wejściowych. Z przeprowadzonej analizy wynika, że poszczególne cechy mają różne zakresy wartości, dlatego proponuje się zastosowanie standaryzacji zgodnie ze wzorem:
\[
x' = \frac{x - \mu}{\sigma},
\]
gdzie $\mu$ oznacza średnią, a $\sigma$ odchylenie standardowe danej cechy.

Parametry standaryzacji powinny zostać wyznaczone wyłącznie na zbiorze treningowym, a następnie wykorzystane również do przeskalowania zbioru testowego.

W analizowanym zbiorze nie stwierdzono braków danych, dlatego nie ma konieczności stosowania metod imputacji. Zidentyfikowane na wykresach obserwacje odstające nie będą usuwane automatycznie, ponieważ zbiór jest stosunkowo mały, a usuwanie takich obserwacji mogłoby prowadzić do utraty istotnych informacji.